% basic
\usepackage[utf8]{inputenc}
\usepackage[t1]{fontenc}
\usepackage{textcomp}
\usepackage{url}
\usepackage{graphicx}
\usepackage{float}
\usepackage{booktabs}
\usepackage{enumitem}
\usepackage{emptypage}
\usepackage{subcaption}
\usepackage{multicol}
\usepackage[usenames,dvipsnames]{xcolor}
\usepackage[polish]{babel}
\usepackage{marginnote}
\usepackage{polski, antpolt} % czcionka antyka półtawskiego
\usepackage{mathpazo}
\usepackage[margin=2.5cm]{geometry} 	     % pakiet potrzebny ustawiania marginesów 
\definecolor{ferrarired}{rgb}{255,40,0}
\usepackage{adjustbox}
\usepackage{tabularx}
\usepackage{booktabs}
\usepackage{amsmath, amssymb, amsfonts, mathtools, amsthm}
\usepackage{mathrsfs}
\usepackage{cancel}
\usepackage{bm}
\newcommand\n{\ensuremath{\mathbb{n}}}
\newcommand\r{\ensuremath{\mathbb{r}}}
\newcommand\z{\ensuremath{\mathbb{z}}}
\renewcommand\o{\ensuremath{\emptyset}}
\newcommand\q{\ensuremath{\mathbb{q}}}
\newcommand\c{\ensuremath{\mathbb{c}}}
\declaremathoperator{\sgn}{sgn}
\declaremathoperator{\fix}{fix}
\usepackage{systeme}
\let\svlim\lim\def\lim{\svlim\limits}
\let\implies\rightarrow
\let\impliedby\leftarrow
\let\iff\leftrightarrow
\let\epsilon\varepsilon
\usepackage{stmaryrd} % for \lightning
\newcommand\contra{\scalebox{1.1}{$\lightning$}}
\let\phi\varphi





% correct
\definecolor{correct}{html}{009900}
\newcommand\correct[2]{\ensuremath{\:}{\color{red}{#1}}\ensuremath{\to }{\color{correct}{#2}}\ensuremath{\:}}
\newcommand\green[1]{{\color{correct}{#1}}}



% horizontal rule
\newcommand\hr{
    \noindent\rule[0.5ex]{\linewidth}{0.5pt}
}


% hide parts
\newcommand\hide[1]{}



% si unitx
\usepackage{siunitx}
\sisetup{locale = fr}
\renewcommand\vec[1]{\mathbf{#1}}
\newcommand\mat[1]{\mathbf{#1}}


% tikz

\usepackage{tikz}
\usepackage{tikz-cd}
\usetikzlibrary{shapes.geometric, fit, intersections, angles, quotes, calc, positioning}
\usetikzlibrary{arrows.meta}
\usepackage{pgfplots}
\pgfplotsset{compat=1.13}
\usetikzlibrary{babel}

\tikzset{
    force/.style={thick, {circle[length=2pt]}-stealth, shorten <=-1pt}
}
\newcommand{\tikzmark}[1]{\tikz[baseline,remember picture] \coordinate (#1) {};}




% theorems
\makeatother
\usepackage{thmtools}
\usepackage[framemethod=tikz]{mdframed}
\mdfsetup{skipabove=1em,skipbelow=0em}


\theoremstyle{definition}

\declaretheoremstyle[
    headfont=\bfseries\sffamily\color{forestgreen!70!black}, bodyfont=\normalfont, postheadspace=\newline,
    mdframed={
        linewidth=1pt,
        rightline=false, topline=false, bottomline=false,
        linecolor=forestgreen, backgroundcolor=forestgreen!3,
    }
]{thmgreenbox}



\declaretheoremstyle[
     headfont=\bfseries\sffamily\color{navyblue!70!black}, bodyfont=\normalfont, postheadspace=\newline,
    mdframed={
        linewidth=1pt,
        rightline=false, topline=false, bottomline=false,
        linecolor=navyblue, backgroundcolor=navyblue!5,
    }
]{thmbluebox}

\declaretheoremstyle[
     headfont=\bfseries\sffamily\color{black}, bodyfont=\normalfont, 
    mdframed={
        linewidth=0.25pt,
	linecolor=black
    }
]{thmblackbox}

\declaretheoremstyle[
headfont=\bfseries\sffamily\color{magenta!70!black}, bodyfont=\normalfont, headpunct = {:},
    mdframed={
        linewidth=1pt,
        rightline=false, topline=false, bottomline=false,
        linecolor=magenta
    }
]{thmmagentaline}

\declaretheoremstyle[
    headfont=\bfseries\sffamily\color{royalblue!70!black}, bodyfont=\normalfont, postheadspace=\newline,
    mdframed={
        linewidth=1pt,
        rightline=false, topline=false, bottomline=false,
        linecolor=royalblue
    }
]{thmblueline}

\declaretheoremstyle[
    headfont=\bfseries\sffamily\color{goldenrod!70!black}, bodyfont=\normalfont, postheadspace=\newline,
    mdframed={
        linewidth=1pt,
        rightline=false, topline=false, bottomline=false,
        linecolor=goldenrod!95!black
    }
]{thmgoldline}

\declaretheoremstyle[
headfont=\bfseries\sffamily\color{blueviolet!70!black}, bodyfont=\normalfont, headpunct = {:}, mdframed={
        linewidth=1pt,
        rightline=false, topline=false, bottomline=false,
        linecolor=blueviolet
    }
]{thmcyanline}

\declaretheoremstyle[
headfont=\bfseries\sffamily\color{orchid!70!black}, bodyfont=\normalfont, headpunct = {:}, mdframed={
        linewidth=1pt,
        rightline=false, topline=false, bottomline=false,
        linecolor=orchid
    }
]{thmcntxtline}

\declaretheoremstyle[
    headfont=\bfseries\sffamily\color{rawsienna!70!black}, bodyfont=\normalfont,postheadspace=\newline, mdframed={
        linewidth=1pt,
        rightline=false, topline=false, bottomline=false,
        linecolor=rawsienna, backgroundcolor=rawsienna!5,
    }
]{thmredbox}

%% goal styles
\declaretheoremstyle[
headfont=\bfseries\sffamily\color{ferrarired!70!black}, bodyfont=\normalfont, headpunct = {:}, mdframed={
        linewidth=1pt,
        rightline=false, topline=false, bottomline=false,
        linecolor=ferrarired
    }
]{thmgoalline}

\declaretheoremstyle[
headfont=\bfseries\sffamily\color{ferrarired!70!black}, bodyfont=\normalfont, headpunct = {:}, mdframed={
        linewidth=1pt,
        rightline=false, topline=false, bottomline=false,
        linecolor=ferrarired, backgroundcolor=ferrarired!5,
    }
]{thmgoalbox}

\declaretheoremstyle[
    headfont=\bfseries\sffamily\color{magenta!70!black}, bodyfont=\normalfont,postheadspace=\newline, mdframed={
        linewidth=1pt,
        rightline=false, topline=false, bottomline=false,
        linecolor=magenta, backgroundcolor=magenta!5,
    }
]{thmmgntbox}

\declaretheoremstyle[
headfont=\bfseries\sffamily\color{orchid!70!black}, bodyfont=\normalfont, headpunct={:}, mdframed={
        linewidth=1pt,
        rightline=false, topline=false, bottomline=false,
        linecolor=orchid, backgroundcolor=orchid!5,
    }
]{thmcntxtbox}

\declaretheoremstyle[
    headfont=\bfseries\sffamily\color{magenta!70!black}, bodyfont=\normalfont, postheadspace=\newline, numbered=no,mdframed={
        linewidth=1pt,
        rightline=false, topline=false, bottomline=false,
        linecolor=magenta, backgroundcolor=magenta!1,
    },
    qed=\qedsymbol
]{thmproofbox}

\declaretheoremstyle[
    headfont=\bfseries\sffamily\color{navyblue!70!black}, bodyfont=\normalfont, postheadspace=\newline,
    numbered=no,
    mdframed={
        linewidth=1pt,
        rightline=false, topline=false, bottomline=false,
        linecolor=navyblue, backgroundcolor=navyblue!3,
    },
]{thmexplanationbox}

\declaretheorem[style=thmgreenbox, name=definicja]{definition}
\declaretheorem[style=thmgoalline, numbered=no, name=cel]{goall}
\declaretheorem[style=thmgoalbox, numbered=no, name=cel]{goal}
\declaretheorem[style=thmcyanline, numbered=no, name=przykład]{eg}
\declaretheorem[style=thmmagentaline, numbered=no, name=pytanie]{question}
\declaretheorem[style=thmmagentaline, numbered=no, name=podsumowanie]{summary}
\declaretheorem[style=thmcntxtbox, numbered=no, name=kontekst]{context}
\declaretheorem[style=thmgoldline, numbered=no, name=intuicja]{intuition}
\declaretheorem[style=thmmgntbox, name=fakt]{fact}
\declaretheorem[style=thmmgntbox, name=twierdzenie]{theorem}
\declaretheorem[style=thmmgntbox, name=lemat]{lemma}
\declaretheorem[style=thmmgntbox, numbered=no, name=wniosek]{corollary}
\declaretheorem[style=thmblackbox, name=ćwiczenie]{ex}

\declaretheorem[style=thmproofbox, name=dowód]{replacementproof}
\renewenvironment{proof}[1][\proofname]{\vspace{-10pt}\begin{replacementproof}}{\end{replacementproof}}


\declaretheorem[style=thmexplanationbox, name=wyjaśnienie]{proposition}
\newenvironment{explanation}[1][]{\vspace{-10pt}\begin{tmpexplanation}}{\end{tmpexplanation}}

\declaretheorem[style=thmblueline, numbered=no, name=spostrzeżenie]{remark}
\declaretheorem[style=thmblueline, numbered=no, name=note]{note}

\newtheorem*{uovt}{uovt}
\newtheorem*{notation}{notation}
\newtheorem*{previouslyseen}{as previously seen}
\newtheorem*{problem}{zadanie}
\newtheorem*{observe}{observe}
\newtheorem*{property}{property}
\newtheorem*{solution}{rozwiązanie}

% figure support
\usepackage{import}
\usepackage{xifthen}
\pdfminorversion=7
\usepackage{pdfpages}
\usepackage{transparent}
\newcommand{\incfig}[1]{%
    \def\svgwidth{\columnwidth}
    \import{./figures/}{#1.pdf_tex}
}


% %http://tex.stackexchange.com/questions/76273/multiple-pdfs-with-page-group-included-in-a-single-page-warning
\pdfsuppresswarningpagegroup=1

\declaremathoperator{\id}{id}
\declaremathoperator{\cl}{cl}
\declaremathoperator{\im}{im}
\declaremathoperator{\aut}{aut}
\declaremathoperator{\bor}{bor}
\declaremathoperator{\def}{def.}
\declaremathoperator{\tor}{tor}
\newcommand{\rotatein}{\rotatebox[origin=c]{90}{$\in$}} %obrocony \in 90stopni counterclockwie

\usepackage{faktor}
\makeatletter
\declarerobustcommand*{\mfaktor}[3][]
{
   { \mathpalette{\mfaktor@impl@}{{#1}{#2}{#3}} }
}
\newcommand*{\mfaktor@impl@}[2]{\mfaktor@impl#1#2}
\newcommand*{\mfaktor@impl}[4]{
   \settoheight{\faktor@zaehlerhoehe}{\ensuremath{#1#2{#3}}}%
   \settoheight{\faktor@nennerhoehe}{\ensuremath{#1#2{#4}}}%
      \raisebox{0.5\faktor@zaehlerhoehe}{\ensuremath{#1#2{#3}}}%
      \mkern-4mu\diagup\mkern-5mu%
      \raisebox{-0.5\faktor@nennerhoehe}{\ensuremath{#1#2{#4}}}%
}
\makeatother
